% Generated draft: requires mathematical and source review.
\documentclass[11pt]{article}
\usepackage[margin=25mm]{geometry}
\usepackage{fontspec}
\setmainfont{Latin Modern Roman}
\usepackage{amsmath,amssymb,mathtools,unicode-math}
\setmathfont{Latin Modern Math}
\usepackage{enumitem,needspace,xurl,hyperref}
\setlength{\parindent}{0pt}
\setlength{\parskip}{5pt}
\setlength{\emergencystretch}{3em}
\newcommand{\N}{\mathbb{N}}
\newcommand{\Z}{\mathbb{Z}}
\newcommand{\Q}{\mathbb{Q}}
\newcommand{\R}{\mathbb{R}}
\newcommand{\C}{\mathbb{C}}
\newcommand{\F}{\mathbb{F}}
\newcommand{\A}{\mathbb{A}}
\newcommand{\PP}{\mathbb{P}}
\newcommand{\E}{\mathbb{E}}
\newcommand{\eps}{\varepsilon}
\newcommand{\dd}{\,\mathrm{d}}
\begin{document}
\section*{UnsolvedMath ID: 800155}
\textbf{Source problem code:} catalogue entry
\subsection{Definitions and mathematical statement}
Let \(\zeta(s)\) denote the Riemann zeta function, defined for \(\operatorname{Re}s>1\) by \(\zeta(s)=\sum_{n=1}^{\infty}n^{-s}\) and elsewhere by meromorphic continuation. Let \(T>0\), let \(0<c_1<c_2\), and let \(\omega\) be a smooth, compactly supported function on \(\R\) satisfying
\[
\operatorname{supp}\omega\subseteq[c_1T,c_2T].
\]
Define its weighted fourth moment by
\[
M_4(T;\omega)=\int_{-\infty}^{\infty}\omega(t)\left|\zeta\!\left(\tfrac12+it\right)\right|^4\,\dd t.
\]
The supplied source asks for the estimate
\[
M_4(T;\omega)=O\!\left(T^{\frac12+\theta+\varepsilon}\right),
\]
where it describes \(\theta\) only as an upper bound in Ramanujan's conjecture.

\textbf{Unresolved qualifications.} The usual asymptotic reading would concern \(T\to\infty\) and every fixed \(\varepsilon>0\), but these quantifiers and the dependence of the implied constant are not specified in the supplied statement. Since the support interval depends on \(T\), an asymptotic formulation also requires a family of weights or a rule for choosing \(\omega\); no normalization, derivative bounds, sign condition, or real versus complex codomain is supplied. The relevant version of Ramanujan's conjecture and the precise meaning of its exponent \(\theta\) are unspecified. The displayed bound concerns the entire fourth moment: the source supplies neither a main term nor a subtraction defining an error term. Whether an error-term estimate was intended remains unresolved, and no such correction is made here.

The source refers to \([\mathrm{Ber},\ \text{Corollary }2]\) as proving an analogous result for modular \(L\)-functions; the supplied record does not identify that reference sufficiently to determine its hypotheses or its bearing on the requested estimate.
\subsection{Short English statement}
The supplied problem asks for a bound of size \(T^{1/2+\theta+\varepsilon}\) for a smoothly weighted fourth moment of the Riemann zeta function. The weight assumptions, the Ramanujan exponent, and whether the intended target is the full moment or an error term require clarification. This is a statement-only draft; the registry's open label is supplied metadata, not an independently verified assessment.
\subsection{Sources}
\begin{itemize}
\item Local UnsolvedMath registry, permanent ID \(800155\), exact source problem code \texttt{catalogue entry}; title: ``Moments of L-functions --- Moments of the Riemann zeta function --- Problem 4.'' The supplied record attributes the problem to Hamieh, Kadiri, and Ng, with proposed year 2024, and labels it open. These are supplied attributions and status metadata, not independently verified conclusions.
\item Hamieh, Kadiri, and Ng, \emph{Moments of L-functions Problem List} (2024), Problem 4. Supplied source URL: \url{https://arxiv.org/abs/2405.17800}. The record reports access on 2026-10-16 and extraction by \texttt{arxiv\_latex}; those provenance details are reproduced without claiming independent source access or verification.
\item The statement cites \([\mathrm{Ber},\ \text{Corollary }2]\) for an analogous result for modular \(L\)-functions. No author identification, title, publication details, or URL for this reference is supplied.
\end{itemize}
\end{document}
